\documentclass[11pt,a4paper]{amsart}
\usepackage{kpfonts}
\usepackage[T1]{fontenc}
\usepackage{mathtools,amssymb,amsthm,amsmath,amsfonts}
\usepackage[square,numbers]{natbib}
\usepackage[dvipsnames]{xcolor}
\usepackage{geometry}
\usepackage{hyperref}
\hypersetup{colorlinks=true,linkcolor=NavyBlue,urlcolor=RubineRed,citecolor=blue}
\geometry{a4paper,total={170mm,257mm},left=20mm,top=20mm}
\numberwithin{equation}{section}
\setcounter{tocdepth}{1}

\title{Probability and Statistical Models in Association Football}
\author{Joyal Liju Jacob}
\date{16 September 2026}

\begin{document}
	
\begin{abstract}
This report examines probability and statistical models that is used to study association football. While the main subject of this report is the paper titled 'A Bayesian Bivariate Conditional Poisson Regression for Goal Dependence in the English Premier League' by Nolan.et al, further models that have been used to model football have also been mentioned. 
\end{abstract}

\maketitle
\tableofcontents

\section{Problem statement and motivation}
Football is the most known sports in the world and it provides a natural setting for applying probability because goals, shots, cards and other events are uncertain, sparse, and generated by two teams that interact strategically. A useful model should assign probabilities to scorelines or events, quantify uncertainty, incorporate team strength and home advantage, and remain interpretable enough to answer questions about match state.

Let $Y_H$ and $Y_A$ be home and away goals respectively. This leads to several questions such as what is the joint distribution
$P(Y_H=y_H,Y_A=y_A)$, how do factors such as team strength and venue affect it, and to what extent are the two scoring processes related? The same framework can be extended to other questions, including shots on target, player position, event timing, league position, and situations in which the two teams have different incentives.

The starting point for this project is Nolan et al.~\cite{nolan2026}, who use a Bayesian bivariate conditional Poisson regression to model dependence between home and away goal counts. The sequence of models below shows how the problem develops from a simple independent Poisson model to a model that can represent positive or negative dependence, followed by several possible directions for further study.

\noindent\rule{\textwidth}{0.8pt}

\section{Nolan et al.'s model development}
This section follows the order of models used by Nolan et al.~\cite{nolan2026}. The starting point is the simplest model, the Independent Poisson Regression.

\subsection*{Independent Poisson regression: the baseline}
Maher~\cite{maher1982} and Dixon and Coles~\cite{dixoncoles1997} use the following baseline idea.

\subsubsection*{Random variable.} For match $m$, define $Y_{Hm},Y_{Am}\in\{0,1,2,\ldots\}$ as home and away goals.
\subsubsection*{Distribution.} The baseline assumes : 

\begin{equation}
Y_{Hm}\sim\operatorname{Poisson}(\lambda_{Hm}),\qquad Y_{Am}\sim\operatorname{Poisson}(\lambda_{Am}),\qquad Y_{Hm}\perp Y_{Am}.
\end{equation}

Thus $P(Y=y)=e^{-\lambda}\lambda^y/y!$ and the scoreline probability is the product of the two marginal probabilities.

\subsubsection*{Football interpretation.} $Y_H$ and $Y_A$ are the final goals scored by the home and away teams.

\subsubsection*{Model parameters.} A common log-link is
\begin{align}
\log\lambda_H&=\mu+\eta+\alpha_{\text{home}}-\delta_{\text{away}}+x_H^{\mathsf T}\beta_H,\\
\log\lambda_A&=\mu+\alpha_{\text{away}}-\delta_{\text{home}}+x_A^{\mathsf T}\beta_A,
\end{align}
where $\mu$ is baseline scoring, $\eta$ is home advantage, $\alpha_i$ is attack strength, $\delta_i$ is defensive strength, and $\beta$ contains covariate effects.

\subsubsection*{Assumptions.} The model assumes integer counts, a constant match-level rate conditional on covariates, equality of the Poisson mean and variance, and conditional independence of the two scores.
\subsubsection*{Problem solved.} It provides a tractable likelihood for scoreline forecasting and estimation of attack, defence and home effects.
\subsubsection*{Limitation addressed.} This is the baseline rather than an extension: it does not address tactical interaction, excess draws, overdispersion, or changing team strength.

\noindent\rule{\textwidth}{0.4pt}

\subsection{Dixon--Coles low-score adjustment}
Dixon and Coles~\cite{dixoncoles1997} retain the Poisson marginals but multiply the joint mass by a correction $\tau(y_H,y_A)$ for the four low scorelines:
\begin{equation}
P(Y_H=y_H,Y_A=y_A)=\tau(y_H,y_A)f_P(y_H;\lambda_H)f_P(y_A;\lambda_A),
\end{equation}
where $\tau$ is a function of dependence parameter $\rho$ for $(0,0),(0,1),(1,0),(1,1)$ and equals one elsewhere.

\subsubsection*{Random variable.} The random variables remain home and away goals, $Y_H$ and $Y_A$.
\subsubsection*{Distribution.} Each variable has a Poisson marginal, but the joint mass is multiplied by the low-score correction $\tau(y_H,y_A)$.
\subsubsection*{Football interpretation.} The variables are the final score, with special treatment of frequent low-scoring outcomes.
\subsubsection*{Model parameters.} Parameters are $\lambda_H,\lambda_A,\rho$, together with attack, defence, home-advantage and covariate parameters.
\subsubsection*{Assumptions.} Conditional on the rates, dependence is concentrated in $(0,0),(0,1),(1,0),(1,1)$ and the correction is one elsewhere.
\subsubsection*{Problem solved.} It improves scoreline and betting-probability forecasting by correcting the low-score region.
\subsubsection*{Limitation addressed.} It addresses independent Poisson misfit for $0$--$0$, $1$--$0$, $0$--$1$ and $1$--$1$, but does not provide a general dependence mechanism for all scorelines.

\noindent\rule{\textwidth}{0.4pt}

\subsection{Bivariate Poisson model}
Karlis and Ntzoufras~\cite{karlis2003} define independent latent counts $U_1\sim\operatorname{Poisson}(\lambda_1)$, $U_2\sim\operatorname{Poisson}(\lambda_2)$, and $U_0\sim\operatorname{Poisson}(\lambda_0)$, then set
\begin{equation}Y_H=U_1+U_0,\qquad Y_A=U_2+U_0.\end{equation}

\subsubsection*{Random variable.} The random variables are home and away goals, constructed from latent counts $U_0,U_1,U_2$.
\subsubsection*{Distribution.} $U_1\sim\operatorname{Poisson}(\lambda_1)$, $U_2\sim\operatorname{Poisson}(\lambda_2)$ and $U_0\sim\operatorname{Poisson}(\lambda_0)$ are independent, with $Y_H=U_1+U_0$ and $Y_A=U_2+U_0$.
\subsubsection*{Football interpretation.} $Y_H$ and $Y_A$ are final goals; $U_0$ represents shared match tempo or common opportunities.
\subsubsection*{Model parameters.} The parameters are $\lambda_0,\lambda_1,\lambda_2$, linked to attack, defence, home advantage and covariates. The marginal means are $\lambda_1+\lambda_0$ and $\lambda_2+\lambda_0$, and covariance is $\lambda_0$.
\subsubsection*{Assumptions.} The latent counts are independent Poisson variables, so dependence is non-negative and each marginal remains equidispersed.
\subsubsection*{Problem solved.} It provides joint scoreline prediction and improves representation of draws.
\subsubsection*{Limitation addressed.} It addresses independence, but cannot express negative goal dependence or flexible dispersion. Nolan et al.~\cite{nolan2026} motivate their conditional-Poisson extension partly because tactical responses can produce negative dependence.

\noindent\rule{\textwidth}{0.4pt}

\subsection{Bayesian bivariate conditional Poisson model (Nolan et al.)}
Nolan et al.~\cite{nolan2026} use a conditional construction in which one goal count is
Poisson and the other goal count is Poisson conditional on the first:
\begin{equation}
Y_1\sim\operatorname{Poisson}(\lambda_1),\qquad
Y_2\mid Y_1=y_1\sim\operatorname{Poisson}\!\left(\lambda_2e^{\theta y_1}\right).
\end{equation}

\subsubsection*{Random variable.} $Y_1$ and $Y_2$ are the home and away goal counts for a match.
\subsubsection*{Distribution.} $Y_1$ has a Poisson marginal distribution, while $Y_2$ has a conditional Poisson distribution given $Y_1$. The resulting joint distribution is a bivariate conditional Poisson distribution.
\subsubsection*{Football interpretation.} The variables represent final home and away goals. The conditional term represents how the scoring of one team changes the expected scoring of the other.
\subsubsection*{Model parameters.} The model has separate marginal means $\lambda_1,\lambda_2$, a dependence parameter $\theta$, regression coefficients for covariates such as attendance and fouls, and Bayesian prior/posterior distributions for the unknown parameters. The log means can include separate home and away effects.
\subsubsection*{Assumptions.} Conditional on the covariates and the first goal count, the second count is Poisson. The regression link is correctly specified, observations are appropriately treated as conditionally independent across matches, and the prior distributions adequately represent parameter uncertainty.
\subsubsection*{Problem solved.} The BCP model produces a joint predictive distribution for home and away goals while allowing match-level covariates and Bayesian uncertainty quantification. It is therefore suited to posterior predictive checks and realistic scoreline probabilities.
\subsubsection*{Limitation addressed.} It addresses the independent Poisson assumption and the bivariate Poisson restriction to non-negative dependence. In Nolan et al.'s application, the dependence parameter can represent either positive or negative association between the two scoring processes.

\noindent\rule{\textwidth}{0.8pt}

\section{Other models across papers}
These models are relevant alternatives or future directions, but they are kept separate
from the main Nolan et al. development.
\subsection{COM--Poisson and dispersion-flexible count models}
Piancastelli et al.~\cite{piancastelli2023} use a flexible multivariate Conway--Maxwell--Poisson construction for English Premier League data. For one team,
\begin{equation}P(Y=y\mid\lambda,\nu)=\frac{\lambda^y}{(y!)^\nu Z(\lambda,\nu)},\qquad Z(\lambda,\nu)=\sum_{k=0}^{\infty}\frac{\lambda^k}{(k!)^\nu}.
\end{equation}

\subsubsection*{Random variable.} $Y$ is a team's match goal count; the multivariate version jointly models the teams.
\subsubsection*{Distribution.} The distribution is COM--Poisson with mass $\lambda^y/((y!)^\nu Z(\lambda,\nu))$.
\subsubsection*{Football interpretation.} $Y$ represents goals scored by a team in a match.
\subsubsection*{Model parameters.} $\lambda>0$ is an intensity and $\nu>0$ controls dispersion: $\nu=1$ is approximately Poisson, $\nu>1$ permits underdispersion and $\nu<1$ permits overdispersion. Regression coefficients use team strength, venue, attendance and other covariates.
\subsubsection*{Assumptions.} The model assumes a COM--Poisson-shaped count distribution and a specified multivariate dependence structure.
\subsubsection*{Problem solved.} It models goal variance that does not equal the mean and can represent richer dependence.
\subsubsection*{Limitation addressed.} It addresses Poisson equidispersion and restrictive dependence, at the cost of a difficult normalising constant, more parameters and overfitting risk.

\subsection{Dynamic Bayesian Poisson model}
Rue and Salvesen~\cite{ruesalvesen2000} model team attack and defence strengths as time-dependent quantities. A representative formulation is
\begin{equation}Y_{Ht}\mid\lambda_{Ht}\sim\operatorname{Poisson}(\lambda_{Ht}),\quad Y_{At}\mid\lambda_{At}\sim\operatorname{Poisson}(\lambda_{At}),\end{equation}
with $\log\lambda_{Ht}=\mu+\eta+a_{H,t}-d_{A,t}$ and $\log\lambda_{At}=\mu+a_{A,t}-d_{H,t}$, while latent strengths evolve as $a_{i,t}=a_{i,t-1}+\epsilon_{i,t}$.

\subsubsection*{Random variable.} The random variables are match goal counts $Y_{Ht},Y_{At}$.
\subsubsection*{Distribution.} Conditional distributions are Poisson.
\subsubsection*{Football interpretation.} The latent states represent changing team form or ability through a season.
\subsubsection*{Model parameters.} Parameters include baseline scoring, home advantage, latent attack and defence strengths, and state-evolution variances.
\subsubsection*{Assumptions.} Conditional on current strengths, matches are conditionally independent and the strength process carries information forward.
\subsubsection*{Problem solved.} It forecasts matches when teams improve, decline, suffer injuries or change form.
\subsubsection*{Limitation addressed.} It addresses constant team parameters in static Poisson regression, but introduces state uncertainty, identifiability issues and computational cost.

\subsection{Ordered probit/logit models for match result}
Goddard~\cite{goddard2005} compares a goal-based bivariate Poisson model with a direct result-based ordered probit model. Let $R\in\{0,1,2\}$ denote away win, draw and home win, and define
\begin{equation}R^*=x^{\mathsf T}\beta+\varepsilon,\qquad\varepsilon\sim N(0,1),\end{equation}
with thresholds $\kappa_1<\kappa_2$ determining the category.

\subsubsection*{Random variable.} The random variable is the categorical match result $R$, not the exact score.
\subsubsection*{Distribution.} $R$ has an ordered categorical distribution: $P(R=0)=\Phi(\kappa_1-x^\top\beta)$, $P(R=1)=\Phi(\kappa_2-x^\top\beta)-\Phi(\kappa_1-x^\top\beta)$, and the remainder is $P(R=2)$.
\subsubsection*{Football interpretation.} $R\in\{0,1,2\}$ denotes away win, draw and home win.
\subsubsection*{Model parameters.} Parameters are $\beta$ and thresholds; covariates may include team ratings, home advantage and form.
\subsubsection*{Assumptions.} The model assumes one latent continuous performance difference and ordered thresholds.
\subsubsection*{Problem solved.} It forecasts win/draw/loss directly when exact score is not the target.
\subsubsection*{Limitation addressed.} It addresses score-model complexity and rare exact scorelines, but loses goal information and imposes an ordinal structure on result categories.

\subsection{Copula-based bivariate goal models}
McHale and Scarf~\cite{mchalescarf2007} join chosen marginal distributions with a copula. If $F_H,F_A$ are marginal CDFs and $C_\theta$ is a copula,
\begin{equation}P(Y_H\le y_H,Y_A\le y_A)=C_\theta\{F_H(y_H),F_A(y_A)\}.
\end{equation}

\subsubsection*{Random variable.} The random variables are home and away goals.
\subsubsection*{Distribution.} Marginals may be discrete Poisson-like goal distributions, joined by a copula $C_\theta$.
\subsubsection*{Football interpretation.} The variables represent final goals and the copula represents association between opponents.
\subsubsection*{Model parameters.} Parameters include marginal rates, team covariate coefficients and copula dependence parameter(s) $\theta$.
\subsubsection*{Assumptions.} The selected marginals and copula are correctly specified.
\subsubsection*{Problem solved.} It provides flexible joint scoreline forecasting.
\subsubsection*{Limitation addressed.} It addresses the bivariate Poisson's particular dependence structure, but discrete copulas are harder to interpret and estimate.

\subsection{Elo/rating plus categorical regression}
Hvattum and Arntzen~\cite{hvattum2010} use Elo ratings as a compact measure of historical team strength. A simple categorical formulation is
\begin{equation}P(R=r)=\frac{\exp(\alpha_r+\beta_r\Delta E+\gamma_r H)}{\sum_{s\in\{0,1,2\}}\exp(\alpha_s+\beta_s\Delta E+\gamma_s H)},
\end{equation}
where $\Delta E$ is the pre-match Elo difference and $H$ is home advantage.

\subsubsection*{Random variable.} $R$ is the win/draw/loss result.
\subsubsection*{Distribution.} Its distribution is multinomial, or ordered categorical in regression variants.
\subsubsection*{Football interpretation.} $R$ records the match result; Elo is a latent state summarising past results and opposition strength.
\subsubsection*{Model parameters.} Parameters are intercepts, regression coefficients, the Elo update constant and optional home/margin adjustments.
\subsubsection*{Assumptions.} Past results can be compressed into a rating and the rating-to-outcome mapping is stable.
\subsubsection*{Problem solved.} It gives fast pre-match forecasting with few covariates.
\subsubsection*{Limitation addressed.} It addresses static team-strength parameters, but does not naturally provide a full scoreline, goal timing or shot-level explanation; ratings can lag sudden changes.

\subsection{Shot-level and time-to-event models}

\subsubsection{Expected goals as logistic regression}
Expected-goals research models the outcome of each shot. Fairchild et al.~\cite{fairchild2018} describe logistic regression of shot success; later work studies player and position corrections~\cite{fairchild2023}. Let $G_j$ indicate whether shot $j$ becomes a goal:
\begin{equation}G_j\mid x_j\sim\operatorname{Bernoulli}(p_j),\qquad\operatorname{logit}(p_j)=x_j^{\mathsf T}\beta.
\end{equation}

\subsubsection*{Random variable.} $G_j\in\{0,1\}$ is the outcome of shot $j$.
\subsubsection*{Distribution.} $G_j$ has a Bernoulli distribution with success probability $p_j$.
\subsubsection*{Football interpretation.} $G_j$ records whether a shot becomes a goal; $x_j$ can contain distance, angle, shot type, assist type, location, player and position.
\subsubsection*{Model parameters.} Parameters are coefficients $\beta$, with optional player/team random effects. Team expected goals is $X_T=\sum_{j\in T}p_j$.
\subsubsection*{Assumptions.} Shots are conditionally independent given features and the log-odds is described by the covariates.
\subsubsection*{Problem solved.} The model measures chance quality and separates shot volume from shot quality.
\subsubsection*{Limitation addressed.} It addresses the goal-count model's failure to distinguish one high-quality chance from many poor shots, but does not by itself model event sequence, defensive pressure or final score.

\subsection{Weibull accelerated failure-time model for the next goal}
Clegg, Song and Cartlidge~\cite{clegg2026} model time to the next goal in-play using a Weibull accelerated failure-time model. For a positive waiting time $T$,
\begin{equation}\log T=\mu+x^{\mathsf T}\beta+\sigma\varepsilon,\qquad S(t)=\Pr(T>t)=\exp[-(t/\lambda)^k].\end{equation}

\subsubsection*{Random variable.} $T$ is the time from the current match state until the next goal, possibly censored at the final whistle.
\subsubsection*{Distribution.} $T$ is Weibull conditional on covariates, with survival function $S(t)=\Pr(T>t)=\exp[-(t/\lambda)^k]$.
\subsubsection*{Football interpretation.} $T$ represents the waiting time until the next goal in an in-play match.
\subsubsection*{Model parameters.} Parameters are shape $k$, baseline scale/intercept, regression coefficients and team-strength effects; covariates can include current score, half, red cards, post-shot expected goals and market information.
\subsubsection*{Assumptions.} The model assumes a specified survival/hazard form and correctly handled censoring.
\subsubsection*{Problem solved.} It forecasts event timing and gives $\Pr(T\le h)=1-S(h)$ for the next $h$ minutes.
\subsubsection*{Limitation addressed.} It addresses the static Poisson model's failure to represent waiting times and evolving match state, but a Weibull hazard may be too restrictive after goals, cards, substitutions or tactical transitions.

\noindent\rule{\textwidth}{0.8pt}

\section{Comparison and research direction}
The models considered here answer different questions. Count models describe goals and scorelines, result models describe win/draw/loss, shot models describe the quality of individual chances, survival models describe when an event occurs, and dynamic models allow team strength to change over time. These are different levels of the same football process rather than direct substitutes for one another.

A possible next step is to compare independent Poisson, Dixon--Coles, bivariate Poisson, COM--Poisson, and negative-binomial models using out-of-sample predictive performance and calibration. The project could then move from final scores to individual events by modelling shots on target, player position, and team or player effects. Predicted shot probabilities could be incorporated into a goal model, while score, cards, substitutions, and elapsed time could be treated as components of a changing match state.

A further question is whether the scoring rate changes when the teams face different incentives, for example when one team needs a win and the other does not. League-table position, match importance, scoreline, and time remaining could then be used to investigate this question.


\newpage


\begin{thebibliography}{99}
\bibitem[Fairchild et al.(2018)]{fairchild2018} Fairchild, A., Pelechrinis, K. and Kokkodis, M. (2018). Spatial analysis of shots in MLS: a model for expected goals and fractal dimensionality. \emph{Journal of Sports Analytics}. \url{https://doi.org/10.3233/JSA-170207}.
\bibitem[Fairchild et al.(2023)]{fairchild2023} Fairchild, A. et al. (2023). A machine learning approach for player and position adjusted expected goals in football (soccer). \emph{Finance Research Letters}, 55, 100034. \url{https://doi.org/10.48550/arXiv.2301.13052}.
\bibitem[Goddard(2005)]{goddard2005} Goddard, J. (2005). Regression models for forecasting goals and match results in association football. \emph{International Journal of Forecasting}, 21, 331--340. \url{https://doi.org/10.1016/j.ijforecast.2004.08.002}.
\bibitem[Hvattum and Arntzen(2010)]{hvattum2010} Hvattum, L. M. and Arntzen, H. (2010). Using ELO ratings for match result prediction in association football. \emph{International Journal of Forecasting}, 26, 460--470. \url{https://doi.org/10.1016/j.ijforecast.2009.10.002}.
\bibitem[Karlis and Ntzoufras(2003)]{karlis2003} Karlis, D. and Ntzoufras, I. (2003). Analysis of sports data by using bivariate Poisson models. \emph{The Statistician}, 52, 381--393. \url{https://doi.org/10.1111/1467-9884.00366}.
\bibitem[Maher(1982)]{maher1982} Maher, M. J. (1982). Modelling association football scores. \emph{Statistica Neerlandica}, 36, 109--118. \url{https://doi.org/10.1111/j.1467-9574.1982.tb00782.x}.
\bibitem[McHale and Scarf(2007)]{mchalescarf2007} McHale, I. G. and Scarf, P. (2007). Modelling soccer matches using bivariate discrete distributions with general dependence structure. \emph{Statistical Modelling}, 7, 147--161. \url{https://www.researchgate.net/publication/4791355_Modelling_soccer_matches_using_bivariate_discrete_distributions_with_general_dependence_structure}.
\bibitem[Nolan et al.(2026)]{nolan2026} Nolan, M., Barreto-Souza, W., Piancastelli, L. S. C. and Sundararajan, R. R. (2026). A Bayesian bivariate conditional Poisson regression for goal dependence in the English Premier League. arXiv:2608.07168. \url{https://arxiv.org/abs/2608.07168}.
\bibitem[Piancastelli et al.(2023)]{piancastelli2023} Piancastelli, L., Friel, N., Barreto-Souza, W. and Ombao, H. (2023). Multivariate Conway--Maxwell--Poisson distribution: Sarmanov method and doubly-intractable Bayesian inference. \emph{Journal of Computational and Graphical Statistics}, 32, 483--500. \url{https://doi.org/10.48550/arXiv.2107.07561}.
\bibitem[Rue and Salvesen(2000)]{ruesalvesen2000} Rue, H. and Salvesen, O. (2000). Prediction and retrospective analysis of soccer matches in a league. \emph{The Statistician}, 49, 399--418.
\bibitem[Dixon and Coles(1997)]{dixoncoles1997} Dixon, M. J. and Coles, S. G. (1997). Modelling association football scores and inefficiencies in the football betting market. \emph{Journal of the Royal Statistical Society: Series C}, 46, 265--280. \url{https://doi.org/10.1111/1467-9876.00065}.
\bibitem[Clegg et al.(2026)]{clegg2026} Clegg, L., Song, Z. and Cartlidge, J. (2026). A market-calibrated accelerated failure time model for in-play football forecasting. arXiv:2605.16066. \url{https://arxiv.org/abs/2605.16066}.
\end{thebibliography}
\end{document}
